\section{Criteris i mètriques d'avaluació}

Una comparació útil necessita definir cada instant del cicle de vida. Si
$a_i$ és l'arribada de la tasca $i$, $s_i$ l'inici i $c_i$ la finalització, el
temps d'espera és $s_i-a_i$ i el temps total al sistema és $c_i-a_i$. La durada
del servei és $c_i-s_i$ en el model sense interrupció. Aquesta definició evita
usar indistintament espera, resposta i execució.

La mitjana resumeix el comportament general, però és sensible a valors extrems
i pot ocultar que una classe concreta queda perjudicada. Per aquest motiu convé
afegir la mediana, el percentil 95 i el màxim del temps d'espera, així com el
nombre de tasques que superen un llindar. Els resultats s'han de calcular també
per prioritat o tipus. Aquesta desagregació connecta amb la literatura sobre
equitat, que estudia com una política pot millorar una mesura global mentre
reparteix de manera desigual el retard \parencite{wierman2011fairness}.

El throughput és el nombre de tasques completades per unitat de temps simulat.
La utilització d'un worker és la proporció del temps observat durant la qual es
troba ocupat. Buzen relaciona throughput, resposta i utilització mitjançant
lleis operacionals de rendiment informàtic \parencite{buzen1976laws}. Aquestes
relacions són útils com a comprovacions de coherència, però no expliquen per si
soles les diferències entre classes.

La longitud de la cua al llarg del temps ajuda a interpretar pics i períodes de
saturació. Un únic valor final no permet distingir una acumulació breu d'una
cua que creix de manera persistent. En escenaris finits també es pot mesurar el
makespan, el temps fins que finalitza l'última tasca. Dues polítiques poden tenir
un makespan semblant i distribuir l'espera de manera molt diferent.

La selecció mínima proposada inclou mitjana, mediana, percentil 95 i màxim
d'espera; temps total al sistema; tasques completades; throughput; utilització;
longitud de cua; i resultats per prioritat. Si s'incorporen fallades, s'afegiran
intents i reintents. Les mètriques addicionals només entraran si responen una
pregunta concreta, perquè massa indicadors poden dificultar la interpretació.
